\section{Generación del fichero fstab}

% Nota: las particiones ya deben estar montadas desde la sección de particionado.
% Si se ha seguido el paso anterior, /home ya está montado en /mnt/home.

\subsection{Generación de {\ttfamily fstab}}

\begin{tcolorbox}[title=Generación de fstab]
    \begin{lstlisting}[style=bashstyle]
genfstab -U /mnt >> /mnt/etc/fstab
    \end{lstlisting}
\end{tcolorbox}
    
\begin{tcolorbox}[title=Comprobamos que se ha generado correctamente]
    \begin{lstlisting}[style=bashstyle]
cat /mnt/etc/fstab
    \end{lstlisting}
\end{tcolorbox}

\subsection{Entrada a la raíz}

  \begin{tcolorbox}[title=Y entramos en la partición {\ttfamily /mnt}]
    \begin{lstlisting}[style=bashstyle]
arch-chroot /mnt
    \end{lstlisting}
  \end{tcolorbox}
\newpage
